% Source: Weinberg GR, physical PDF pp. 360--363
%         (printed pp. 338--341).
% Coverage: Section 11.8, Comoving Coordinates; equations
%           (11.8.1)--(11.8.19).
% Figures/tables/footnotes: no figures or tables; reference callout 31.
% Status: source-reviewed and compile-clean.
% Uncertainties: the factor-of-two source anomaly following case (B) is
%                documented at the correction below; the curvature conversion
%                is documented at Eqs. (11.8.14)--(11.8.17).

\subsection{Comoving Coordinates}\label{sec:11.8}

As a further preparation for our treatment of gravitational collapse, and
also to lay a groundwork for our discussion of cosmology in Chapter~14, we
now construct a very useful set of coordinates, the \emph{comoving coordinate
system},\hyperref[ch11-ref-31]{\textsuperscript{31}} which incorporates a more
natural separation between space and time than that provided by the standard
coordinates used in the last section.

Imagine a finite region of space filled with a dense cloud of freely falling
particles. Each particle is assumed to carry along a little clock, and is
given a fixed set of spatial coordinates, which can be defined as the
coordinates \(x^i\) of the particle, in some arbitrary system, when its own
clock reads \(t=0\). (The rules for setting these different clocks are
discussed below.) The spacetime coordinates \(x^\mu\) of any event are defined
by taking \(\mathbf{x}\) as the spatial coordinate label of the particle that
is just going by when and where the event occurs, and by taking \(t\) as the
time then shown on that particle's clock. We may think of the coordinate mesh
as being dragged along by the cloud of particles, with time defined by clocks
stuck on the mesh. This coordinate system will be useful throughout the
region occupied by the particle cloud, for whatever interval of time in which
particle trajectories do not cross.

The metric in comoving coordinates is characterized by certain specially
simple features. First, we note that the clocks are in free fall and therefore
tell proper time, so the proper-time interval between two points
\((\mathbf{x},t)\) and \((\mathbf{x},t+\dd t)\) on a given particle's trajectory
is just \(\dd t\). In the binding signature this says
\[
  \dd\tau^2=-g_{\mu\nu}\,\dd x^\mu \dd x^\nu=-g_{tt}\,\dd t^2,
\]
and therefore
\begin{equation}
  g_{tt}=-1.
  \tag{11.8.1}\label{eq:11.8.1}
\end{equation}
Also, we note that the particle trajectory
\(\mathbf{x}=\text{constant}\), \(t=\tau\), satisfies the equation of free
fall, so
\[
  \frac{\dd^2x^i}{\dd\tau^2}
  +\Gamma^i{}_{\mu\nu}
   \frac{\dd x^\mu}{\dd\tau}\frac{\dd x^\nu}{\dd\tau}
  =0.
\]
Using Eq.~\eqref{eq:11.8.1}, this gives
\[
  0=\Gamma^i{}_{tt}
  =\frac12 g^{ij}
   \left(2\partial_tg_{jt}-\partial_jg_{tt}\right)
  =g^{ij}\partial_tg_{jt},
\]
or, since \(g^{ij}\) is generally a nonsingular matrix,
\begin{equation}
  \partial_tg_{it}=0.
  \tag{11.8.2}\label{eq:11.8.2}
\end{equation}

We have kept open the option of setting the clocks attached to the different
particles in an arbitrary fashion. Suppose that we reset these clocks by a
transformation
\begin{equation}
  t'=t+f(\mathbf{x}),\qquad x'^i=x^i.
  \tag{11.8.3}\label{eq:11.8.3}
\end{equation}
The new metric will have the elements
\begin{align}
  g'_{tt}&=-1,
  \tag{11.8.4}\label{eq:11.8.4}\\
  g'_{ti}&=g_{ti}+\partial_i f,
  \tag{11.8.5}\label{eq:11.8.5}\\
  g'_{ij}
  &=g_{ij}
    -g_{ti}\partial_jf-g_{tj}\partial_if
    -\partial_if\,\partial_jf.
  \tag{11.8.6}\label{eq:11.8.6}
\end{align}
It would be a great simplification if the function \(f\) could be chosen so
that the two terms in Eq.~\eqref{eq:11.8.5} cancel, giving \(g'_{ti}=0\).
There are two important cases where this is possible:

\emph{(A)} Suppose that we can reset all clocks so that all particles are at
rest at a time \(t=0\). This assumption can be given an absolute physical
significance by interpreting it to mean that, for each particle \(P\) at
\(t=0\), it is possible to find a locally inertial coordinate system
\(\xi^\mu\) in which the separation between \(P\) and neighboring particles is
purely spatial,
\[
  \left.
  \frac{\partial \xi^0}{\partial x^i}
  \right|_{t=0,\,\mathbf{x}=\mathbf{x}_P}
  =0,
\]
and in which the movement of \(P\) in a time interval \(\dd t\) is purely
temporal,
\[
  \left.
  \frac{\dd\xi^i}{\dd t}
  \right|_{t=0,\,\mathbf{x}=\mathbf{x}_P}
  =0.
\]
The metric in this locally inertial system is the Minkowski metric
\(\eta_{\mu\nu}\), so the time--space components of the metric in the
comoving system at \(t=0\) are
\[
  g_{ti}(\mathbf{x}_P,0)
  =
  \eta_{\mu\nu}
  \frac{\partial\xi^\mu}{\partial t}
  \frac{\partial\xi^\nu}{\partial x^i}
  =0.
\]
With Eq.~\eqref{eq:11.8.2}, it follows that \(g_{ti}\) vanishes everywhere,
so the metric is given by
\begin{equation}
  \dd s^2=-\dd t^2+g_{ij}(\mathbf{x},t)\,\dd x^i \dd x^j.
  \tag{11.8.7}\label{eq:11.8.7}
\end{equation}

\emph{(B)} If the metric is manifestly spherically symmetric, then the line
element must have the general form with which we started in the last section,
that is,
\[
  \begin{aligned}
  \dd s^2={}&-C(r,t)\,\dd t^2+D(r,t)\,\dd r^2+2E(r,t)\,\dd r\,\dd t\\
         &+F(r,t)r^2(\dd\theta^2+\sin^2\theta\,\dd\varphi^2).
  \end{aligned}
\]
The only nonvanishing time--space component \(g_{ti}\) is
\(g_{tr}=E\), and Eq.~\eqref{eq:11.8.2} then tells us that \(E\) is
time-independent, so
\[
  g_{tr}=E(r),\qquad g_{t\theta}=g_{t\varphi}=0.
\]
We can therefore eliminate the components \(g_{ti}\) by resetting the clocks
as in Eq.~\eqref{eq:11.8.3}, with
\[
  f(r)=-\int^r E(\tilde r)\,\dd\tilde r.
\]
% VERIFY SOURCE: The printed text states g_tr=2E and f=-2\int E\,dr,
% even though the displayed line element contains 2E\,dr\,dt and hence
% g_tr=E.  The two factors of two have been corrected for tensor consistency.
Using Eq.~\eqref{eq:11.8.4} and dropping primes, the metric is now of the form
\begin{equation}
  \dd s^2
  =
  -\dd t^2+U(r,t)\,\dd r^2
  +V(r,t)(\dd\theta^2+\sin^2\theta\,\dd\varphi^2),
  \tag{11.8.8}\label{eq:11.8.8}
\end{equation}
with \(U\) and \(V\) new unknown functions that replace \(D\) and \(F\).

It is of course possible to construct coordinate systems of this sort even
if the cloud of freely falling particles is purely imaginary. In differential
geometry, coordinate systems satisfying Eqs.~\eqref{eq:11.8.1} and
\eqref{eq:11.8.2} are called \emph{Gaussian}, and if \(g_{ti}\) vanishes, so
that the line element takes the form \eqref{eq:11.8.7}, then we call the
coordinates \emph{Gaussian normal}. However, these coordinate systems find
their most important applications to systems that actually do consist of a
freely falling fluid. In this case the fluid velocity four-vector by
definition has zero space component,
\begin{equation}
  U^i=0,
  \tag{11.8.9}\label{eq:11.8.9}
\end{equation}
and since \(U^\mu\) is normalized so that
\begin{equation}
  g_{\mu\nu}U^\mu U^\nu=-1,
  \tag{11.8.10}\label{eq:11.8.10}
\end{equation}
[see Eq.~\eqref{eq:5.4.4}], the time component of \(U^\mu\) must be
\begin{equation}
  U^t=(-g_{tt})^{-1/2}=1.
  \tag{11.8.11}\label{eq:11.8.11}
\end{equation}

We shall be working only with spherically symmetric comoving coordinate
systems, with line element \eqref{eq:11.8.8}. The nonvanishing elements of
the metric tensor and its inverse are
\begin{equation}
  \begin{gathered}
  g_{rr}=U,\qquad
  g_{\theta\theta}=V,\qquad
  g_{\varphi\varphi}=V\sin^2\theta,\qquad
  g_{tt}=-1,\\
  g^{rr}=U^{-1},\qquad
  g^{\theta\theta}=V^{-1},\qquad
  g^{\varphi\varphi}=(V\sin^2\theta)^{-1},\qquad
  g^{tt}=-1.
  \end{gathered}
  \tag{11.8.12}\label{eq:11.8.12}
\end{equation}
The nonvanishing elements of the affine connection are readily calculated as
\begin{equation}
  \begin{gathered}
  \Gamma^r{}_{rr}=\frac{U'}{2U},\qquad
  \Gamma^r{}_{\theta\theta}=-\frac{V'}{2U},\qquad
  \Gamma^r{}_{\varphi\varphi}=-\frac{V'}{2U}\sin^2\theta,\qquad
  \Gamma^r{}_{rt}=\Gamma^r{}_{tr}=\frac{\dot U}{2U},\\
  \Gamma^\theta{}_{r\theta}=\Gamma^\theta{}_{\theta r}
    =\frac{V'}{2V},\qquad
  \Gamma^\theta{}_{t\theta}=\Gamma^\theta{}_{\theta t}
    =\frac{\dot V}{2V},\qquad
  \Gamma^\theta{}_{\varphi\varphi}=-\sin\theta\cos\theta,\\
  \Gamma^\varphi{}_{r\varphi}=\Gamma^\varphi{}_{\varphi r}
    =\frac{V'}{2V},\qquad
  \Gamma^\varphi{}_{t\varphi}=\Gamma^\varphi{}_{\varphi t}
    =\frac{\dot V}{2V},\qquad
  \Gamma^\varphi{}_{\theta\varphi}
    =\Gamma^\varphi{}_{\varphi\theta}=\cot\theta,\\
  \Gamma^t{}_{rr}=\frac{\dot U}{2},\qquad
  \Gamma^t{}_{\theta\theta}=\frac{\dot V}{2},\qquad
  \Gamma^t{}_{\varphi\varphi}=\frac{\dot V}{2}\sin^2\theta.
  \end{gathered}
  \tag{11.8.13}\label{eq:11.8.13}
\end{equation}
(A prime or dot denotes \(\partial_r\) or \(\partial_t\), respectively.) From
the definition of the Ricci tensor in the binding curvature convention we
obtain the independent nonzero components:
% MODERNIZATION CHECK: every component in (11.8.14)--(11.8.17) is the
% negative of Weinberg's printed Ricci component.
\begin{align}
  R_{rr}
  ={}&
  -\frac{V''}{V}
  +\frac{V'^2}{2V^2}
  +\frac{U'V'}{2UV}
  +\frac{\ddot U}{2}
  -\frac{\dot U^2}{4U}
  +\frac{\dot U\dot V}{2V},
  \tag{11.8.14}\label{eq:11.8.14}\\
  R_{\theta\theta}
  ={}&
  1-\frac{V''}{2U}
  +\frac{V'U'}{4U^2}
  +\frac{\ddot V}{2}
  +\frac{\dot V\dot U}{4U},
  \tag{11.8.15}\label{eq:11.8.15}\\
  R_{tt}
  ={}&
  -\frac{\ddot U}{2U}
  -\frac{\ddot V}{V}
  +\frac{\dot U^2}{4U^2}
  +\frac{\dot V^2}{2V^2},
  \tag{11.8.16}\label{eq:11.8.16}\\
  R_{tr}
  ={}&
  -\frac{\dot V'}{V}
  +\frac{V'\dot V}{2V^2}
  +\frac{\dot U V'}{2UV}.
  \tag{11.8.17}\label{eq:11.8.17}
\end{align}
Also, it again follows from the spherical symmetry of the metric that
\begin{equation}
  R_{\varphi\varphi}=R_{\theta\theta}\sin^2\theta,
  \tag{11.8.18}\label{eq:11.8.18}
\end{equation}
and
\begin{equation}
  R_{r\theta}=R_{r\varphi}=R_{\theta\varphi}
  =R_{\theta t}=R_{\varphi t}=0.
  \tag{11.8.19}\label{eq:11.8.19}
\end{equation}
